% GENERATED FILE. Edit canonical lessons, then run npm run generate:books.
% canonical-source-hash: fnv1a64:a2358c5885063c0e
% canonical-lessons: TE-C393-prayojanam, TE-C393-labhapadu, TE-C393-nastapovu, TE-C393-hanikaramaina, TE-C393-upayogapadu

\chapter{Gains and Losses}
\label{ch:te-gains-and-losses}
\hlchaptermodality{393}

\begin{chapteropening}
\textbf{By the end of this chapter:} \emph{I can say who gains, who loses and what helps or harms.}

Complete the last lesson of chapter 393: I can say who gains, who loses and what helps or harms.
\end{chapteropening}

\section[\texorpdfstring{\te{ప్రయోజనం} (prayōjanaṁ)}{ప్రయోజనం (prayōjanaṁ)}]{\te{ప్రయోజనం} (prayōjanaṁ) — a benefit; a purpose}

\label{lesson:TE-C393-prayojanam}



\begin{quote}
Before the new one: say the Telugu for to arise, then the Telugu for to crop up.
\end{quote}

\subsection*{You'll want to know: \te{ప్రయోజనం}}
\textbf{\te{ప్రయోజనం}} — \emph{prayōjanaṁ} — \textquotedblleft{}a benefit; a purpose\textquotedblright{}.

\begin{quote}
\te{ఆ} \te{కొత్త} \te{రోడ్డు} \te{వల్ల} \te{ఎవరికి} \te{ప్రయోజనం}?
\end{quote}

\begin{quote}
\emph{ā kotta rōḍḍu valla evariki prayōjanaṁ?}
\end{quote}

\begin{quote}
Who benefits from that new road?
\end{quote}

Plural: \te{ప్రయోజనాలు}, \emph{prayōjanālu}.

\te{ప్రయోజనం} \te{లేదు}, it is no use.

\begin{cousinweb}
\textbf{\te{ప్రయోజనం}} (\emph{prayōjanaṁ}) is borrowed from Sanskrit \emph{prayojana}, \textquotedblleft{}occasion, motive, purpose, aim\textquotedblright{} (Monier-Williams); in Telugu it is also the good a thing does.
\end{cousinweb}

\begin{grammarlens}[title={What you've built}]
One more word for giving reasons and opinions.
\end{grammarlens}

\subsection*{Your turn}
\begin{itemize}
\raggedright
  \item \emph{Say it:} \emph{prayōjanaṁ}
  \item \emph{Say it:} \emph{prayōjanaṁ}, once more
  \item \emph{Say it:} \emph{ā kotta rōḍḍu valla evariki prayōjanaṁ?} — Who benefits from that new road?
  \item \emph{Recall:} say the Telugu for to arise, then the Telugu for to crop up, then say \emph{prayōjanaṁ} again
\end{itemize}

\begin{tcolorbox}[breakable,colback=teal!4,colframe=teal!35!black,title={Before you move on}]
What does \te{ప్రయోజనం} mean? (\textquotedblleft{}A benefit; a purpose\textquotedblright{}.) Say it once more.
\end{tcolorbox}
\section[\texorpdfstring{\te{లాభపడు} (lābhapaḍu)}{లాభపడు (lābhapaḍu)}]{\te{లాభపడు} (lābhapaḍu) — to benefit; to gain}

\label{lesson:TE-C393-labhapadu}



\begin{quote}
Before the new one: say the Telugu for to crop up, then the Telugu for a benefit.
\end{quote}

\subsection*{You'll want to know: \te{లాభపడు}}
\textbf{\te{లాభపడు}} — \emph{lābhapaḍu} — \textquotedblleft{}to benefit; to gain\textquotedblright{}.

\begin{quote}
\te{ఆ} \te{వ్యాపారంలో} \te{మా} \te{మామయ్య} \te{బాగా} \te{లాభపడ్డాడు}.
\end{quote}

\begin{quote}
\emph{ā vyāpāraṁlō mā māmayya bāgā lābhapaḍḍāḍu.}
\end{quote}

\begin{quote}
My uncle gained a lot from that business.
\end{quote}

Past (I …): \te{లాభపడ్డాను}, \emph{lābhapaḍḍānu}.

\te{లాభం}, profit, which you know, + \te{పడు}: to come into profit.

\begin{cousinweb}
\textbf{\te{లాభపడు}} (\emph{lābhapaḍu}) is built from \te{లాభం}, from Sanskrit \emph{lābha}, \textquotedblleft{}gain\textquotedblright{} (CDIAL 11018), and the native verb \emph{paḍu}, \textquotedblleft{}to feel, to undergo\textquotedblright{} (DEDR 3853).
\end{cousinweb}

\begin{grammarlens}[title={What you've built}]
One more word for giving reasons and opinions.
\end{grammarlens}

\subsection*{Your turn}
\begin{itemize}
\raggedright
  \item \emph{Say it:} \emph{lābhapaḍu}
  \item \emph{Say it:} \emph{lābhapaḍu}, once more
  \item \emph{Say it:} \emph{ā vyāpāraṁlō mā māmayya bāgā lābhapaḍḍāḍu.} — My uncle gained a lot from that business.
  \item \emph{Recall:} say the Telugu for to crop up, then the Telugu for a benefit, then say \emph{lābhapaḍu} again
\end{itemize}

\begin{tcolorbox}[breakable,colback=teal!4,colframe=teal!35!black,title={Before you move on}]
What does \te{లాభపడు} mean? (\textquotedblleft{}To benefit; to gain\textquotedblright{}.) Say it once more.
\end{tcolorbox}
\section[\texorpdfstring{\te{నష్టపోవు} (naṣṭapōvu)}{నష్టపోవు (naṣṭapōvu)}]{\te{నష్టపోవు} (naṣṭapōvu) — to suffer a loss}

\label{lesson:TE-C393-nastapovu}



\begin{quote}
Before the new one: say the Telugu for a benefit, then the Telugu for to benefit.
\end{quote}

\subsection*{You'll want to know: \te{నష్టపోవు}}
\textbf{\te{నష్టపోవు}} — \emph{naṣṭapōvu} — \textquotedblleft{}to suffer a loss\textquotedblright{}.

\begin{quote}
\te{వర్షాల} \te{వల్ల} \te{రైతులు} \te{చాలా} \te{నష్టపోయారు}.
\end{quote}

\begin{quote}
\emph{varṣāla valla raitulu cālā naṣṭapōyāru.}
\end{quote}

\begin{quote}
The farmers lost a great deal because of the rains.
\end{quote}

Past (I …): \te{నష్టపోయాను}, \emph{naṣṭapōyānu}.

\te{నష్టం}, a loss, which you know, + \te{పోవు}, to go.

\begin{cousinweb}
\textbf{\te{నష్టపోవు}} (\emph{naṣṭapōvu}) is built from \te{నష్టం}, from Sanskrit \emph{naṣṭa}, \textquotedblleft{}lost, destroyed\textquotedblright{} (CDIAL 7027), and the native verb \emph{pōvu}, \textquotedblleft{}to go, to be lost\textquotedblright{} (DEDR 4572).
\end{cousinweb}

\begin{grammarlens}[title={What you've built}]
One more word for giving reasons and opinions.
\end{grammarlens}

\subsection*{Your turn}
\begin{itemize}
\raggedright
  \item \emph{Say it:} \emph{naṣṭapōvu}
  \item \emph{Say it:} \emph{naṣṭapōvu}, once more
  \item \emph{Say it:} \emph{varṣāla valla raitulu cālā naṣṭapōyāru.} — The farmers lost a great deal because of the rains.
  \item \emph{Recall:} say the Telugu for a benefit, then the Telugu for to benefit, then say \emph{naṣṭapōvu} again
\end{itemize}

\begin{tcolorbox}[breakable,colback=teal!4,colframe=teal!35!black,title={Before you move on}]
What does \te{నష్టపోవు} mean? (\textquotedblleft{}To suffer a loss\textquotedblright{}.) Say it once more.
\end{tcolorbox}
\section[\texorpdfstring{\te{హానికరమైన} (hānikaramaina)}{హానికరమైన (hānikaramaina)}]{\te{హానికరమైన} (hānikaramaina) — harmful}

\label{lesson:TE-C393-hanikaramaina}



\begin{quote}
Before the new one: say the Telugu for to benefit, then the Telugu for to suffer a loss.
\end{quote}

\subsection*{You'll want to know: \te{హానికరమైన}}
\textbf{\te{హానికరమైన}} — \emph{hānikaramaina} — \textquotedblleft{}harmful\textquotedblright{}.

\begin{quote}
\te{పొగ} \te{తాగడం} \te{ఆరోగ్యానికి} \te{హానికరమైన} \te{అలవాటు}.
\end{quote}

\begin{quote}
\emph{poga tāgaḍaṁ ārōgyāniki hānikaramaina alavāṭu.}
\end{quote}

\begin{quote}
Smoking is a habit that harms your health.
\end{quote}

\te{హాని}, harm, + -\te{కర}, causing, + -\te{మైన}.

\begin{cousinweb}
\textbf{\te{హానికరమైన}} (\emph{hānikaramaina}) is built on Sanskrit \emph{hānikara}, \textquotedblleft{}injurious, detrimental\textquotedblright{} (Monier-Williams), from \emph{hāni}, \textquotedblleft{}loss, decrease\textquotedblright{}, with Telugu -\te{మైన} for an adjective.
\end{cousinweb}

\begin{grammarlens}[title={What you've built}]
One more word for giving reasons and opinions.
\end{grammarlens}

\subsection*{Your turn}
\begin{itemize}
\raggedright
  \item \emph{Say it:} \emph{hānikaramaina}
  \item \emph{Say it:} \emph{hānikaramaina}, once more
  \item \emph{Say it:} \emph{poga tāgaḍaṁ ārōgyāniki hānikaramaina alavāṭu.} — Smoking is a habit that harms your health.
  \item \emph{Recall:} say the Telugu for to benefit, then the Telugu for to suffer a loss, then say \emph{hānikaramaina} again
\end{itemize}

\begin{tcolorbox}[breakable,colback=teal!4,colframe=teal!35!black,title={Before you move on}]
What does \te{హానికరమైన} mean? (\textquotedblleft{}Harmful\textquotedblright{}.) Say it once more.
\end{tcolorbox}
\section[\texorpdfstring{\te{ఉపయోగపడు} (upayōgapaḍu)}{ఉపయోగపడు (upayōgapaḍu)}]{\te{ఉపయోగపడు} (upayōgapaḍu) — to be useful; to come in handy}

\label{lesson:TE-C393-upayogapadu}



\begin{quote}
Before the new one: say the Telugu for to suffer a loss, then the Telugu for harmful.
\end{quote}

\subsection*{You'll want to know: \te{ఉపయోగపడు}}
\textbf{\te{ఉపయోగపడు}} — \emph{upayōgapaḍu} — \textquotedblleft{}to be useful; to come in handy\textquotedblright{}.

\begin{quote}
\te{మీ} \te{సలహా} \te{నాకు} \te{చాలా} \te{ఉపయోగపడింది}.
\end{quote}

\begin{quote}
\emph{mī salahā nāku cālā upayōgapaḍindi.}
\end{quote}

\begin{quote}
Your advice was very useful to me.
\end{quote}

Past (it …): \te{ఉపయోగపడింది}, \emph{upayōgapaḍindi}.

The person helped takes -\te{కి}, and the thing that helps is the subject: its past is shown as \te{ఉపయోగపడింది}, it came in handy.

\begin{cousinweb}
\textbf{\te{ఉపయోగపడు}} (\emph{upayōgapaḍu}) is built from Sanskrit \emph{upayoga}, \textquotedblleft{}use, application\textquotedblright{} (Monier-Williams), the word in \te{ఉపయోగించు}, and the native verb \emph{paḍu} (DEDR 3853).
\end{cousinweb}

\begin{grammarlens}[title={What you've built}]
One more word for giving reasons and opinions.
\end{grammarlens}

\subsection*{Your turn}
\begin{itemize}
\raggedright
  \item \emph{Say it:} \emph{upayōgapaḍu}
  \item \emph{Say it:} \emph{upayōgapaḍu}, once more
  \item \emph{Say it:} \emph{mī salahā nāku cālā upayōgapaḍindi.} — Your advice was very useful to me.
  \item \emph{Recall:} say the Telugu for to suffer a loss, then the Telugu for harmful, then say \emph{upayōgapaḍu} again
\end{itemize}

\begin{tcolorbox}[breakable,colback=teal!4,colframe=teal!35!black,title={Before you move on}]
What does \te{ఉపయోగపడు} mean? (\textquotedblleft{}To be useful; to come in handy\textquotedblright{}.) Say it once more.
\end{tcolorbox}
